% Research note: conditional finite-probability theorem.
% The assumption (GSR) has NOT been proved for actual Collatz/Syracuse orbits.
% Do NOT input into main.tex until separately reviewed.
\subsection*{A block-geometric shell-return criterion}

Keep the notation of Section~\ref{sec:tail-programme}, with
$N_{X,j}$ the number of recorded visits to the shell $I_{X,j}$
before first passage. For each fixed $X,j$, suppose that the
logarithmic densities of the events $\{N_{X,j}\ge k\}$ exist, and
write
\[
 p_{X,j}(k):=\mathbb P_{\log}(N_{X,j}\ge k),\qquad k\ge1.
\]
The following statement is strictly conditional; it does not assert
that the actual Syracuse dynamics has geometric return tails.

\begin{proposition}[Block-geometric return criterion]
\label{prop:block-geometric-shell-return}
Suppose there exist an integer $b\ge1$ and $0\le\rho<1$,
independent of all sufficiently large $X$, all shell indices $j$,
and all $k\ge1$, such that
\begin{equation}\label{eq:block-geometric-shell-assumption}
 p_{X,j}(k+b)\le\rho\,p_{X,j}(k).
\end{equation}
Then
\[
 \mathbb E_{\log}[N_{X,j}^{2}]
 \le \frac{b^{2}(1+\rho)}{(1-\rho)^{2}},
\]
uniformly in $X,j$. Consequently, assuming the definition of $T_X$
and the dyadic-shell criterion from Section~\ref{sec:tail-programme},
\[
 T_X \le
 \frac{4b^{2}(1+\rho)}{(1-\rho)^{2}}
 X(\log X)^{-\varepsilon_0/2}
 =o(X).
\]
\end{proposition}
\begin{proof}
For an integer-valued, nonnegative random variable $N$,
\[
 N^{2}=\sum_{k\ge1}(2k-1)\mathbf 1_{\{N\ge k\}},
\qquad
\mathbb E[N^{2}]=\sum_{k\ge1}(2k-1)\mathbb P(N\ge k).
\]
These sums are finite in our first-passage setting for each fixed
$X,j$: before reaching the lower threshold the trajectory cannot
revisit a state, and each shell contains finitely many integers.
By monotonicity and the assumed $b$-step inequality, for
$mb+1\le k\le(m+1)b$,
\[
 p_{X,j}(k)\le p_{X,j}(mb+1)
 \le \rho^{m}p_{X,j}(1)\le\rho^{m}.
\]
Thus
\[
 \begin{aligned}
 \mathbb E_{\log}[N_{X,j}^{2}]
 &\le \sum_{m\ge0}\rho^{m}
    \sum_{k=mb+1}^{(m+1)b}(2k-1)\\
 &= b^{2}\sum_{m\ge0}(2m+1)\rho^{m}
 =\frac{b^{2}(1+\rho)}{(1-\rho)^{2}}.
 \end{aligned}
\]
Apply the dyadic-shell second-moment criterion.
\end{proof}

\paragraph{Equivalent conditional interpretation.}
When $p_{X,j}(k)>0$, assumption
\eqref{eq:block-geometric-shell-assumption} says that, under the
\emph{logarithmic distribution of initial values conditioned on at
least $k$ visits}, the chance of reaching $k+b$ visits is at most
$\rho$. No independent restart of the deterministic orbit is
assumed. The claim must be proved for this induced conditional
measure, and it cannot be read off from the ideal renewal model.

\paragraph{Research obstacle.}
The local-time bounds in the renewal model and finite sampled
histograms do not establish a uniform $\rho<1$ for actual orbits.
The effect of the induced distribution of shell entrance states,
particularly on rare trajectories with many returns, is open.
The separate residual term $R_X$ is unaffected.

\paragraph{The independent valuation model does satisfy a geometric bound.}
The obstruction is the transfer to actual orbits, not the renewal model.
For the idealised walk $S_{n+1}=S_n+\xi_{n+1}$ with independent
$\xi=g\log2-\log3$, $\mathbb P(g=k)=2^{-k}$ for $k\ge1$,
let $N_I$ count all visits to any fixed interval
$I=(a,a+\log2]$, starting from an arbitrary initial point.
Then
\begin{equation}\label{eq:model-geometric-rho}
 \mathbb P(N_I\ge k+1)\le (1-p_0)\mathbb P(N_I\ge k),
 \qquad
 p_0=\frac18\left(1-\sqrt{\frac38}\right)>0.
\end{equation}
Indeed, after a visit at $x\in I$, the independent next increment
satisfies $\xi\ge4\log2-\log3$ with probability $1/8$.
On that event the walk rises at least
$\delta=\log(8/3)$ above the upper endpoint of $I$.
Moreover
\[
 \mathbb E[e^{-\xi/2}]
 =\frac{\sqrt3}{2^{3/2}-1}<1.
\]
Thus the exponential process $e^{-S_n/2}$ is a nonnegative
supermartingale. The standard maximal inequality for such processes
bounds the probability of ever falling $\delta$ below the new
position by $e^{-\delta/2}=\sqrt{3/8}$. Consequently, at
\emph{each} visit the conditional probability of no further visit
is at least $p_0$; the strong Markov property gives
\eqref{eq:model-geometric-rho}. In particular,
\[
 \sup_{a,x}\mathbb E_x[N_I^2]
 \le\frac{2-p_0}{p_0^2}<832.
\]
This is a deliberately crude explicit bound for the independent
model. It does \emph{not} yield a return-tail bound for deterministic
Collatz orbits, whose valuation sequence has no such independently
resampled future under the induced starting-value measure.
